\hsection{Chapter 3}
\sol{304} $P_2(x)=P_1(x)+\frac12f''(a)(x-a)^2$: the quadratic correction has the sign of $f''(a)$. For $e^x$ ($f''>0$) it adds to $P_1$: the linear value $1.5$ was too small; for $\log(1+x)$ and $\sqrt[3]x$ ($f''<0$) it subtracts: the linear values were too large. This is the precise form of the hint in 302.

\sol{305} $P_3'(x)=f'(a)+f''(a)(x-a)+\frac12f'''(a)(x-a)^2$, $P_3''(x)=f''(a)+f'''(a)(x-a)$, $P_3'''(x)=f'''(a)$. At $x=a$ all terms with a factor $(x-a)$ vanish: $P_3(a)=f(a)$, $P_3'(a)=f'(a)$, $P_3''(a)=f''(a)$, $P_3'''(a)=f'''(a)$.

\sol{306} The $k$-th derivative of $(x-a)^m$ is $0$ at $x=a$ unless $m=k$, when it is $k!$ (for $m<k$ the term is differentiated away; for $m>k$ a factor $(x-a)^{m-k}$ remains). So $P_n^{(k)}(a)=\frac{f^{(k)}(a)}{k!}\cdot k!=f^{(k)}(a)$ for $0\le k\le n$.

\sol{307} All derivatives of $e^x$ at $0$ are $1$: $P_9=\sum_{k=0}^9\frac{x^k}{k!}=1+x+\frac{x^2}2+\frac{x^3}6+\frac{x^4}{24}+\frac{x^5}{120}+\frac{x^6}{720}+\frac{x^7}{5040}+\frac{x^8}{40320}+\frac{x^9}{362880}$. $\sin$: derivatives cycle $\cos,-\sin,-\cos,\sin$, at $0$: $1,0,-1,0,\dots$: $x-\frac{x^3}{3!}+\frac{x^5}{5!}-\frac{x^7}{7!}+\frac{x^9}{9!}$. $\cos$: $1-\frac{x^2}{2!}+\frac{x^4}{4!}-\frac{x^6}{6!}+\frac{x^8}{8!}$. $\log(1+x)$: $f^{(k)}(0)=(-1)^{k-1}(k-1)!$ (313a), so $c_k=\frac{(-1)^{k-1}}k$: $x-\frac{x^2}2+\frac{x^3}3-\frac{x^4}4+\frac{x^5}5-\frac{x^6}6+\frac{x^7}7-\frac{x^8}8+\frac{x^9}9$.

\sol{310} (a) $f(0)=4$, $f'(x)=6x^2-6x$ so $f'(0)=0$, $f''(x)=12x-6$ so $f''(0)=-6$, $f'''=12$, higher derivatives $0$. $P_0=P_1=4$; $P_2=4-3x^2$; $P_n=4-3x^2+2x^3=f$ for all $n\ge3$. (b) Yes. $f$ and $P_{n,a}$ are both polynomials of degree $\le n$ (expand $f$ in powers of $x-a$: $f(x)=\sum_{k=0}^nc_k(x-a)^k$); by Oefening 306 both have the same derivatives of order $0,\dots,n$ at $a$, and the derivatives at $a$ determine the coefficients $c_k=f^{(k)}(a)/k!$. So $f=P_{n,a}$ exactly.

\sol{311} If $f$ is even then $f'$ is odd, $f''$ even, etc.; an odd function (defined at $0$) vanishes at $0$, so all odd-order derivatives of an even $f$ vanish at $0$: its Taylor approximations at $0$ contain only even powers ($\cos$). For odd $f$ only odd powers occur ($\sin$, $\tan$, $\log\frac{1+x}{1-x}$).

\sol{312} With $T=\tan$, $T'=1/\cos^2$, $T''=2\sin/\cos^3$, $T'''=(2\cos^2+6\sin^2)/\cos^4=(2+4\sin^2)/\cos^4$, and so on; at $0$: $0,1,0,2,0,16$. Result $x+\frac13x^3+\frac2{15}x^5$; the organised way is Oefening 317 (worked in the text).

\sol{313} (a) $\frac{d^k}{dx^k}\log(1+x)=\frac{(-1)^{k-1}(k-1)!}{(1+x)^k}$ (each differentiation multiplies by $-k/(1+x)$), so at $0$: $(-1)^{k-1}(k-1)!$, i.e.\ $0!,-1!,2!,-3!,\dots$ For $\log(1-x)$ the chain rule gives an extra factor $-1$ each time: $\frac{d^k}{dx^k}\log(1-x)=-\frac{(k-1)!}{(1-x)^k}$, at $0$: $-(k-1)!$, all negative. (b) $c_k=-\frac1k$: $\log(1-x)\approx-x-\frac{x^2}2-\frac{x^3}3-\frac{x^4}4-\frac{x^5}5$. (c) Substituting $-x$ in $x-\frac{x^2}2+\frac{x^3}3-\frac{x^4}4+\frac{x^5}5$ gives $-x-\frac{x^2}2-\frac{x^3}3-\frac{x^4}4-\frac{x^5}5$: the same (even powers keep their sign, odd powers flip).

\sol{315} Substitute $u=-x^2$ in $\sqrt{1+u}=1+\frac12u-\frac18u^2+\dots$: $\sqrt{1-x^2}\approx1-\frac12x^2-\frac18x^4$. At $x=\frac12$: $\sqrt{1-\frac14}=\frac12\sqrt3\approx1-\frac18-\frac1{128}=\frac{128-16-1}{128}=\frac{111}{128}$ ($=0.8672$; true value $0.8660$).

\sol{318} $u=-\frac12x^2+\frac1{24}x^4+\text{h.o.t.}$ has lowest order $2$, so $u^2=\frac14x^4+\text{h.o.t.}$ has lowest order $4$ and $u^3,u^4,\dots$ have order $\ge6$. In $\log(1+u)=u-\frac12u^2+\frac13u^3-\dots$ everything from $u^3$ on is beyond order $4$, and within $u^2$ only $\frac14x^4$ matters. So the terms kept, $u-\frac12u^2$ with $u$ up to $x^4$, are exactly all contributions of order $\le4$.

\sol{321} $u=x-\frac16x^3$ tends to $0$, so substitute in $e^u=1+u+\frac{u^2}2+\frac{u^3}6+\frac{u^4}{24}+\frac{u^5}{120}$. Up to order $5$: $u^2=x^2-\frac13x^4$, $u^3=x^3-\frac12x^5$, $u^4=x^4$, $u^5=x^5$. Sum: $1+x-\frac16x^3+\frac12x^2-\frac16x^4+\frac16x^3-\frac1{12}x^5+\frac1{24}x^4+\frac1{120}x^5=1+x+\frac12x^2-\frac18x^4-\frac3{40}x^5$. (Since $x-\frac16x^3$ is $\sin x$ up to order $4$, this agrees with Oefening 320a up to $x^4$.)

\sol{322} (a) $\log\frac12=\log(1+x)$ with $x=-\frac12$: $P_n(-\frac12)=-\sum_{k=1}^n\frac{1}{k\,2^k}$: $-0.5$, $-0.625$, $-0.6667$, $-0.6823$, $-0.6885$, $-0.6911$, $-0.6922$, $-0.6927$, $-0.6929$, $-0.6931$ (for $n=1,\dots,10$) against $\log\frac12=-0.693147$: degree $10$ is needed for the first three decimals ($0.693$). (b) $\log0=\log(1+(-1))$ would need $x=-1$: $P_n(-1)=-\sum_{k=1}^n\frac1k$, the harmonic sums, which go to $-\infty$ very slowly; no degree gives a usable number (and $\log0$ does not exist). (c) $\log1.9$: $x=0.9$ is close to $1$, so the terms $\frac{0.9^k}k$ shrink slowly and many dozens of terms are needed for three decimals. $\log2.1$: $x=1.1>1$, the terms $\frac{1.1^k}k$ grow, the approximations never settle: it does not work at all. (d) $\log2=-\log\frac12\approx0.693$. (e) With $\log\frac{1+x}{1-x}=2(x+\frac{x^3}3+\frac{x^5}5+\dots)$ and $x=\frac13$ (so that $\frac{1+x}{1-x}=2$): $2(0.33333+0.01235+0.00082+0.00007)=0.69313$, three decimals right after four terms. Faster because $x=\frac13$ is small (each term shrinks by a factor $9$) and only odd powers occur.

\sol{324} $\dfrac x{\tan x}=\dfrac{x\cos x}{\sin x}=\cos x\cdot\dfrac1{\sin x/x}$. As $x\to0$: $\cos x\to1$ and $\frac{\sin x}x\to1$ (standard limit), so by the product and quotient rules the limit is $1\cdot\frac11=1$.

\sol{327} For $x\to0^+$ put $x=\frac1u$ with $u\to\infty$: $x\log x=\frac1u\log\frac1u=-\frac{\log u}u\to0$ by 326(a). (Also with l'H\^opital on $\frac{\log x}{1/x}$: $\frac{1/x}{-1/x^2}=-x\to0$.)

\sol{328} $\frac{x^k}{k!}$ for $k=0,1,2,\dots$ gives $1,x,\frac{x^2}2,\frac{x^3}6,\dots$; $(-1)^k\frac{x^{2k+1}}{(2k+1)!}$ gives $x,-\frac{x^3}{6},\frac{x^5}{120},\dots$; $(-1)^k\frac{x^{2k}}{(2k)!}$ gives $1,-\frac{x^2}2,\frac{x^4}{24},\dots$; $(-1)^{k+1}\frac{x^k}k$ gives $x,-\frac{x^2}2,\frac{x^3}3,\dots$: exactly the Taylor polynomials of Oefening 307.

\sol{329} $e^{ix}=\sum_k\frac{(ix)^k}{k!}$ and $i^k$ cycles through $1,i,-1,-i$. The even $k=2m$ give $i^{2m}=(-1)^m$: $\sum_m(-1)^m\frac{x^{2m}}{(2m)!}=\cos x$; the odd $k=2m+1$ give $i^{2m+1}=i(-1)^m$: $i\sum_m(-1)^m\frac{x^{2m+1}}{(2m+1)!}=i\sin x$. So $e^{ix}=\cos x+i\sin x$, Definition 1.8.

\sol{330} (a) $(1+x)(1-x+x^2-x^3+\dots)=1-x+x^2-x^3+\dots+x-x^2+x^3-\dots=1$: all terms cancel except $1$. Dividing by $1+x$: $\frac1{1+x}=1-x+x^2-x^3+\dots$ (b) $x=0$: $1=1$, fine. $x=-\frac12$: $\frac1{1/2}=2$ and $1+\frac12+\frac14+\dots=2$, fine. $x=-1$: $\frac10$ against $1+1+1+\dots$: both ``infinite'', not a meaningful equation. $x=1$: $\frac12$ against $1-1+1-1+\dots$, whose partial sums are $1,0,1,0,\dots$: nonsense. (c) $f(x)=(1+x)^{-1}$ has $f^{(k)}(x)=(-1)^kk!(1+x)^{-k-1}$, so $f^{(k)}(0)/k!=(-1)^k$: the Taylor series is $1-x+x^2-x^3+\dots$, the same as in (a). (d) $\frac d{dx}\big(x-\frac{x^2}2+\frac{x^3}3-\frac{x^4}4+\dots\big)=1-x+x^2-x^3+\dots$, which is the series of $\frac1{1+x}=\frac d{dx}\log(1+x)$: differentiating term by term gives the right derivative (for $|x|<1$, where the series converge).
